\scanpage{8}
qui donne la description suivante de (13) comme extension déduite de l'extension évidente de \deleted{\unsure{par}} $\mathbb Z/2\mathbb Z$ par $\mathbb Z$, grâce à l'hom diag. $\mathbb Z\;-\;\mathbb Z^I$

$(15)$\quad
\diagram{\begin{tikzcd}[ampersand replacement=\&, column sep=large,row sep=large]
0\arrow[r,no head]\&\mathbb Z\arrow[r,no head,"2"]\arrow[d]\&\mathbb Z\arrow[r]\arrow[d]\&\mathbb Z/2\mathbb Z\arrow[r]\arrow[d,"\unsure{\wr|2}"]\&0\\
0\arrow[r,no head]\&\mathbb Z^I\arrow[r]\&\mathcal E_I\arrow[r]\&\mu_2\arrow[r,no head]\&0
\end{tikzcd}}

Donc pour un $\mu_2$-torseur $\varepsilon$ i.e. un ens à deux éléments $\varepsilon$, on a pour $x,y\in\varepsilon$ un isom canonique de $\mathbb Z^I$-torseurs
\[
(16)\qquad\underbrace{\operatorname{Hom}}_{\mathcal E_I(\varepsilon)}(x,y)
=T_{x-y}\wedge_{\mathbb Z}\mathbb Z^I\simeq
\begin{cases}
\mathbb Z^I&\text{si }x=y\\
\overset{\deleted{\unsure{\mathbb Z}}}{\mathrm{Imp}}^I&\text{si }x\ne y
\end{cases}
\]
\margin{où Imp est l'ens des entiers impairs. \unsure{dans}}
où \deleted{$\unsure{T_{xy}}$} $T_{x-y}$ désigne le $\mathbb Z$-torseur image inverse de $x-y\in\mathbb Z/2\mathbb Z$ dans $\mathbb Z$ (i.e. formé des entiers pairs si $x=y$, des entiers impairs si $x\ne y$), et où l'extension du groupe d'opérateurs des torseurs est \unsure{relatif} à l'hom. diagonal $\mathbb Z\;-\;\mathbb Z^I$.

Ceci s'applique en particulier au cas de $\varepsilon=\omega(I)$, et achève de donner la description algébrique complète de $S_0\mathcal T_{0,I}$, en termes de $I\in\operatorname{Ob}\mathrm{Ens}_3$.

Cette description est compatible avec les opérations de $\unsure{\mathfrak S}_I$. On trouvera donc un
\mspage{7}
